\section{Methodology}

\subsection{Formality Measurement}

We measure formality using the DeBERTa Large Formality Ranker \citep{deberta2024formality}, a transformer model fine-tuned on the GYAFC dataset with 87.8\% classification accuracy. The model produces continuous scores in $[-1, 1]$, where higher values indicate greater formality.

For each conversation turn, we extract the text and compute its formality score. We then compute the \textbf{formality delta} between human input and AI response:

\begin{equation}
\Delta_{\text{AI}\to\text{H}} = |f_{\text{AI}_1} - f_{\text{H}_1}|
\end{equation}

Lower delta indicates stronger accommodation---the AI matches the human's formality level more closely.

\subsection{Dataset and Preprocessing}

We use the Anthropic hh-rlhf dataset \citep{anthropic2022hh}, which contains 170,000+ human-AI conversations with clear turn structure. We apply the following filters:

\begin{itemize}
    \item \textbf{Minimum turns}: $\geq$2 turns per participant
    \item \textbf{Language}: English only
    \item \textbf{Message length}: $\geq$5 tokens
\end{itemize}

After filtering, we retain 111,039 conversations for primary analysis.

\subsection{Analysis Framework}

\textbf{H-E1 (Existence)}: We compute Bidirectional Convergence Scores and test for non-trivial variance (SD $> 0.15$).

\textbf{H-M1 (User Adaptation)}: We compute lagged cross-correlation between AI formality and subsequent human formality.

\textbf{H-M2 (AI Accommodation)}: We test whether AI formality correlates with human formality in turn-1 pairs (success: $|r| > 0.1$, $p < 0.001$).

\textbf{H-M3 (Engagement Link)}: We divide conversations into terciles by formality delta magnitude and measure continuation rates per tercile.
